\documentclass[a4paper,12pt]{article}

\usepackage[T2A]{fontenc}
\usepackage[utf8]{inputenc}
\usepackage[russian]{babel}
\usepackage{amsmath,amssymb,amsthm,mathtools}
\usepackage{helvet}
\renewcommand{\familydefault}{\sfdefault}
\usepackage{icomma}
\usepackage{geometry}
\usepackage{microtype}
\usepackage{tikz}
\usetikzlibrary{calc}
\usepackage{tkz-euclide}
\usepackage{pgfplots}
\pgfplotsset{compat=1.18}
\usepackage{tabularx,booktabs,longtable}
\usepackage{enumitem}
\usepackage{hyperref}
\usepackage{parskip}
\usepackage{fancyhdr}
\usepackage{setspace}
\usepackage{xcolor}
\usepackage{titlesec}

\geometry{margin=2.1cm}
\onehalfspacing

\definecolor{accent}{HTML}{008080}
\definecolor{darkaccent}{HTML}{004D4D}
\definecolor{textgray}{HTML}{333333}
\definecolor{softgray}{HTML}{F2F4F4}
\definecolor{muted}{HTML}{6B7C7C}

\titleformat{\section}
  {\Large\bfseries\color{darkaccent}}
  {}{0pt}{}

\titleformat{\subsection}
  {\large\bfseries\color{accent}}
  {}{0pt}{}

\titlespacing*{\section}{0pt}{0pt}{0.8em}
\titlespacing*{\subsection}{0pt}{1.1em}{0.45em}

\hypersetup{
  colorlinks=true,
  linkcolor=darkaccent,
  urlcolor=accent,
  pdftitle={Чек-лист по планиметрии},
  pdfauthor={Лёвин Артём Александрович}
}

\pagestyle{fancy}
\fancyhf{}
\fancyhead[L]{\small\color{muted}Планиметрия: стратегия решения}
\fancyhead[R]{\small\color{muted}Николь Саркисьянц}
\fancyfoot[C]{\thepage}
\renewcommand{\headrulewidth}{0.3pt}
\setlength{\headheight}{15pt}

\setlist[itemize]{
  leftmargin=1.5em,
  itemsep=0.25em,
  topsep=0.3em
}

\setlist[enumerate]{
  leftmargin=1.8em,
  itemsep=0.55em,
  topsep=0.4em
}

\renewcommand{\arraystretch}{1.25}

\tkzSetUpLine[
  line width=1pt,
  color=accent
]

\tkzSetUpPoint[
  color=accent,
  fill=accent,
  size=3
]

\tkzSetUpLabel[
  color=black,
  font=\small
]

\newcommand{\important}[1]{%
  \textbf{\color{darkaccent}#1}%
}

\newcommand{\checkq}[1]{%
  \textbf{\color{accent}#1}%
}

\begin{document}

% =========================================================
% ТИТУЛЬНЫЙ ЛИСТ
% =========================================================

\begin{titlepage}
\thispagestyle{empty}
\centering

\vspace*{3.2cm}

{\fontsize{19}{23}\selectfont
\bfseries\color{darkaccent}
Чек-лист\par}

\vspace{0.45cm}

{\fontsize{26}{31}\selectfont
\bfseries
Планиметрия:\par}

\vspace{0.15cm}

{\fontsize{20}{25}\selectfont
стратегия решения, высоты, площади, окружность и подобие\par}

\vspace{2.2cm}

{\large ЕГЭ по математике\par}

\vfill

{\large\bfseries
Лёвин Артём Александрович\par}

\vspace{0.2cm}

{\large
эксклюзивно для Николь Саркисьянц\par}

\vspace{0.8cm}

{\large \today\par}

\vspace{2.5cm}

\begin{tikzpicture}[remember picture,overlay]
\draw[
  accent!55,
  line width=1.15pt
]
  ([xshift=1.7cm,yshift=1.8cm]current page.south west)
  .. controls
  ([xshift=5.0cm,yshift=2.55cm]current page.south west)
  and
  ([xshift=-5.3cm,yshift=1.05cm]current page.south east)
  ..
  ([xshift=-1.7cm,yshift=1.75cm]current page.south east);
\end{tikzpicture}

\end{titlepage}

% =========================================================
% 1
% =========================================================

\section{1. Архитектура решения геометрической задачи}

В планиметрии особенно полезна последовательность вопросов,
которая быстро превращает рисунок в систему связей.

\begin{enumerate}

  \item
  \checkq{В какую фигуру входит искомая величина?}

  Угол может быть частью треугольника или четырёхугольника;
  длина --- стороной нескольких фигур;
  площадь --- частью более крупной композиции.

  \item
  \checkq{Есть ли прямоугольные треугольники?}

  Если они уже присутствуют, сразу проверяйте
  теорему Пифагора и тригонометрические отношения.
  Если их нет, рассмотрите полезное проведение
  высоты или перпендикуляра.

  \item
  \checkq{Есть ли возможность получить окружность?}

  Два прямых угла, опирающихся на один отрезок,
  часто дают циклический четырёхугольник.

  \item
  \checkq{Что общего у нескольких фигур?}

  Общая сторона, высота, угол, гипотенуза или дуга
  часто создают короткую связь между величинами.

  \item
  \checkq{Какая теория связывает известное с искомым?}

  Проверяйте теорему Пифагора, синус и косинус,
  площади, подобие, вписанные углы,
  сумму углов четырёхугольника.

  \item
  \checkq{Не стал ли расчёт чрезмерно громоздким?}

  Если быстро появляются новые неизвестные
  и длинные уравнения, полезно проверить
  альтернативный метод.

\end{enumerate}

\important{Практический критерий.}
Хорошее решение обычно использует уже заложенную
в конфигурации структуру и вводит минимум новых величин.

\subsection{Работа с дополнительной переменной}

Если искомую длину удобно обозначить через $x$,
выражайте через неё остальные неизвестные элементы.

Для получения единственного численного ответа
должно хватать независимых связей.
При двух переменных обычно ищут две независимые связи.

\subsection{Чертёж на экзамене}

Чертёж должен сохранять смысл условия:

\begin{itemize}
  \item прямой угол выглядит прямым;
  \item тупой угол выглядит тупым;
  \item высота действительно перпендикулярна
  соответствующей прямой;
  \item равные стороны и углы удобно отмечены.
\end{itemize}

Точная масштабность не требуется.
Главная задача рисунка --- сделать связи видимыми
и снизить риск логической ошибки.

% =========================================================
% 2
% =========================================================

\clearpage
\section{2. Высота и прямоугольные треугольники}

\important{Высота треугольника}
--- перпендикуляр, проведённый из вершины
к прямой, содержащей противоположную сторону.

Поэтому основание высоты может находиться:

\begin{itemize}
  \item внутри противоположной стороны;
  \item в вершине;
  \item за пределами отрезка на продолжении
  соответствующей стороны.
\end{itemize}

Последний случай особенно важен
для тупоугольных треугольников.

\begin{center}
\begin{tikzpicture}[scale=1.05]

  \tkzDefPoint(0,0){B}
  \tkzDefPoint(3,0){C}
  \tkzDefPoint(4,2.4){A}
  \tkzDefPoint(4,0){H}

  \tkzDrawSegments(A,B A,C B,H A,H)
  \tkzDrawSegment[dashed,color=muted](C,H)

  \tkzMarkRightAngle[size=0.28](A,H,C)

  \tkzDrawPoints(A,B,C,H)

  \tkzLabelPoints[below](B,C,H)
  \tkzLabelPoints[above](A)

\end{tikzpicture}
\end{center}

Если $AH$ --- высота к прямой $BC$, площадь вычисляется
по стороне исходного треугольника:

\[
S_{ABC}
=
\frac12\,BC\cdot AH.
\]

Длина вспомогательного отрезка $BH$ или $CH$
не заменяет основание $BC$.

\subsection{Главный набор для прямоугольного треугольника}

Для прямоугольного треугольника
с катетами $a,b$ и гипотенузой $c$:

\[
a^2+b^2=c^2.
\]

Для острого угла $\alpha$:

\[
\sin\alpha
=
\frac{\text{противолежащий катет}}
     {\text{гипотенуза}},
\]

\[
\cos\alpha
=
\frac{\text{прилежащий катет}}
     {\text{гипотенуза}}.
\]

Если известен $\sin\alpha$, то для острого $\alpha$

\[
\cos\alpha
=
\sqrt{1-\sin^2\alpha}.
\]

Знак берётся положительный, поскольку

\[
0^\circ<\alpha<90^\circ.
\]

\important{Контроль результата:}

\[
0\le \sin\alpha\le 1,
\qquad
0\le \cos\alpha\le 1
\]

для острых углов.

Значение вне этого диапазона сигнализирует
об ошибке в вычислениях или в выборе отношения сторон.

% =========================================================
% 3
% =========================================================

\clearpage
\section{3. Один пример --- два метода}

Рассмотрим согласованную конфигурацию:

\[
AC=BC,
\qquad
AB=25,
\qquad
AH=20,
\]

где

\[
AH\perp BC,
\]

а точка $H$ лежит на отрезке $BC$.

Требуется найти $AC$.

\begin{center}
\begin{tikzpicture}[scale=1.05]

  \tkzDefPoint(0,0){B}
  \tkzDefPoint(3,0){H}
  \tkzDefPoint(4.1667,0){C}
  \tkzDefPoint(3,4){A}

  \tkzDrawPolygon(A,B,C)
  \tkzDrawSegment(A,H)

  \tkzMarkRightAngle[size=0.28](A,H,C)

  \tkzDrawPoints(A,B,C,H)

  \tkzLabelPoints[below](B,H,C)
  \tkzLabelPoints[above](A)

  \tkzLabelSegment[left](A,H){$20$}
  \tkzLabelSegment[above left](A,B){$25$}

\end{tikzpicture}
\end{center}

\important{Проверка конфигурации.}
При этих числах треугольник получается остроугольным;
точка $H$ должна располагаться внутри $BC$.

\subsection{Метод 1. Дважды применить теорему Пифагора}

В $\triangle ABH$:

\[
BH
=
\sqrt{AB^2-AH^2}
=
\sqrt{25^2-20^2}
=
\sqrt{225}
=
15.
\]

Пусть

\[
AC=BC=x.
\]

Тогда

\[
CH
=
BC-BH
=
x-15.
\]

В $\triangle ACH$:

\[
x^2
=
20^2+(x-15)^2.
\]

Раскрываем скобки:

\[
x^2
=
400+x^2-30x+225.
\]

Отсюда

\[
30x=625,
\]

\[
\boxed{x=\frac{125}{6}}.
\]

Следовательно,

\[
\boxed{AC=\frac{125}{6}}.
\]

\subsection{Метод 2. Через тригонометрию}

Пусть

\[
\beta=\angle ABC.
\]

В прямоугольном $\triangle ABH$:

\[
\sin\beta
=
\frac{20}{25}
=
\frac45,
\]

\[
\cos\beta
=
\frac{15}{25}
=
\frac35.
\]

Поскольку

\[
AC=BC,
\]

углы при основании равны, поэтому

\[
\angle ACB
=
180^\circ-2\beta.
\]

Следовательно,

\[
\sin\angle ACB
=
\sin(180^\circ-2\beta)
=
\sin2\beta.
\]

По формуле двойного угла:

\[
\sin2\beta
=
2\sin\beta\cos\beta
=
2\cdot\frac45\cdot\frac35
=
\frac{24}{25}.
\]

В прямоугольном $\triangle ACH$:

\[
\sin\angle ACH
=
\frac{AH}{AC}.
\]

Так как

\[
\angle ACH=\angle ACB,
\]

получаем

\[
\frac{24}{25}
=
\frac{20}{AC}.
\]

Отсюда

\[
\boxed{AC=\frac{125}{6}}.
\]

\important{Вывод.}
Оба способа корректны.
Сначала полезно проверять теорему Пифагора.
Тригонометрия становится сильной альтернативой,
когда нужные углы удобно связаны с известными сторонами.

% =========================================================
% 4
% =========================================================

\clearpage
\section{4. Метод равенства площадей}

Если для одного треугольника известны разные основания
и соответствующие высоты, площадь связывает их напрямую:

\[
S
=
\frac12 a h_a
=
\frac12 b h_b.
\]

После сокращения коэффициента $\frac12$ получаем

\[
\boxed{a h_a=b h_b}.
\]

\subsection{Пример}

Пусть две стороны треугольника равны $9$ и $6$.

Высота к стороне $9$ равна $4$.

Найдём высоту $h$ к стороне $6$:

\[
\frac12\cdot9\cdot4
=
\frac12\cdot6\cdot h.
\]

Отсюда

\[
36=6h,
\]

\[
\boxed{h=6}.
\]

\important{Сигнал к применению.}
Если в условии одновременно встречаются несколько высот
одного треугольника, метод равенства площадей
следует проверить одним из первых.

\subsection{Почему этот путь часто выигрывает}

Попытка перейти к нескольким прямоугольным треугольникам
может породить дополнительные отрезки и неизвестные.

Формула площади использует высоты напрямую
и обычно даёт одно короткое уравнение.

% =========================================================
% 5
% =========================================================

\clearpage
\section{5. Искомый угол внутри четырёхугольника}

В остром $\triangle ABC$ проведены высоты $BD$ и $CE$,
пересекающиеся в точке $O$.

Пусть

\[
\angle A=62^\circ.
\]

Требуется найти

\[
\angle DOE.
\]

\begin{center}
\begin{tikzpicture}[scale=1.12]

  \tkzDefPoint(0,0){A}
  \tkzDefPoint(4,0){B}
  \tkzDefPoint(1.643,3.090){C}

  \tkzDefPoint(0.8816,1.6581){D}
  \tkzDefPoint(1.6432,0){E}
  \tkzDefPoint(1.6432,1.2532){O}

  \tkzDrawPolygon(A,B,C)
  \tkzDrawSegments(B,D C,E)

  \tkzMarkRightAngle[size=0.24](B,D,A)
  \tkzMarkRightAngle[size=0.24](C,E,A)

  \tkzMarkAngle[size=0.62](B,A,C)
  \tkzLabelAngle[pos=0.9](B,A,C){$62^\circ$}

  \tkzDrawPoints(A,B,C,D,E,O)

  \tkzLabelPoints[below](A,B,E)
  \tkzLabelPoints[above](C,D)
  \tkzLabelPoint[right](O){$O$}

\end{tikzpicture}
\end{center}

Искомый угол является полноценным углом
четырёхугольника $ADOE$.

В нём

\[
\angle ADO=90^\circ,
\]

\[
\angle AEO=90^\circ,
\]

\[
\angle DAE=62^\circ.
\]

По сумме углов четырёхугольника:

\[
\angle DOE
=
360^\circ
-
90^\circ
-
90^\circ
-
62^\circ.
\]

Поэтому

\[
\boxed{\angle DOE=118^\circ}.
\]

\important{Методическая мысль.}
Перед вычислениями спросите:

\[
\text{«В какую фигуру входит искомый угол?»}
\]

Иногда короткая связь находится в четырёхугольнике
быстрее, чем в цепочке треугольников.

% =========================================================
% 6
% =========================================================

\clearpage
\section{6. Окружность как дополнительная конструкция}

\subsection{Теорема Фалеса и обратный переход}

Если $BC$ --- диаметр окружности,
а точка $D$ лежит на окружности, то

\[
\boxed{\angle BDC=90^\circ}.
\]

Обратная идея особенно полезна.

Если

\[
\angle BDC=90^\circ,
\]

то точка $D$ лежит на окружности
с диаметром $BC$.

Следовательно, если два прямоугольных треугольника
имеют общую гипотенузу $BC$,
их вершины прямых углов лежат
на одной окружности с диаметром $BC$.

\begin{center}
\begin{tikzpicture}[scale=0.95]

  \tkzDefPoint(0,0){O}
  \tkzDefPoint(-3,0){B}
  \tkzDefPoint(3,0){C}
  \tkzDefPoint(1.8,2.4){D}
  \tkzDefPoint(-1.8,2.4){E}

  \tkzDrawCircle(O,B)

  \tkzDrawSegments(
    B,C
    B,D
    D,C
    B,E
    E,C
  )

  \tkzMarkRightAngle[size=0.25](B,D,C)
  \tkzMarkRightAngle[size=0.25](B,E,C)

  \tkzDrawPoints(B,C,D,E,O)

  \tkzLabelPoints[below](B,O,C)
  \tkzLabelPoints[above](D,E)

\end{tikzpicture}
\end{center}

Для четырёх точек $B,C,D,E$ на одной окружности
можно использовать свойства вписанных углов.

Например, углы, опирающиеся на одну хорду $ED$
и лежащие в одном сегменте, равны:

\[
\boxed{\angle EBD=\angle ECD}.
\]

\subsection{Связь с высотами}

В треугольнике $ABC$ пусть $D$ и $E$
--- основания высот из $B$ и $C$.

Тогда

\[
\angle BDC
=
\angle BEC
=
90^\circ.
\]

Поэтому точки

\[
B,\ C,\ D,\ E
\]

лежат на одной окружности
с диаметром $BC$.

Это открывает доступ к равенствам вписанных углов
и часто ведёт к подобию треугольников.

\important{Важная пара дополнительных построений:}
перпендикуляр и окружность.

В сложной конфигурации стоит проверять оба варианта.

% =========================================================
% 7
% =========================================================

\clearpage
\section{7. Подобие и общие элементы}

Подобие особенно ценно,
потому что связывает углы и длины
в разных частях рисунка.

В конфигурации с высотами

\[
BD\perp AC,
\qquad
CE\perp AB
\]

треугольники $ABD$ и $ACE$ имеют:

\[
\angle ADB
=
\angle AEC
=
90^\circ,
\]

\[
\angle BAD
=
\angle CAE
=
\angle A.
\]

Следовательно,

\[
\boxed{\triangle ABD\sim\triangle ACE}
\]

по двум углам.

\subsection{Общая высота и отношение площадей}

Пусть точки $A,K,C$ лежат на одной прямой,
а вершина $B$ общая для треугольников
$ABK$ и $KBC$.

Высота из $B$ к прямой $AC$
одинакова для обоих треугольников.

Поэтому

\[
\frac{S_{ABK}}{S_{KBC}}
=
\frac{AK}{KC}.
\]

Это следствие удобно получать
каждый раз из формулы площади:

\[
\frac{
  \frac12 AK\cdot h
}{
  \frac12 KC\cdot h
}
=
\frac{AK}{KC}.
\]

Общий множитель

\[
\frac12 h
\]

сокращается.

\subsection{Пример с долями}

В $\triangle MJQ$ точка $T$ лежит на $MQ$ и

\[
MT:TQ=5:3,
\]

\[
S_{MJQ}=1.
\]

Общая высота из $J$ даёт

\[
S_{MJT}:S_{TJQ}
=
5:3.
\]

Всего имеется $8$ равных долей площади.

Поэтому

\[
\boxed{S_{MJT}=\frac58},
\]

\[
\boxed{S_{TJQ}=\frac38}.
\]

\subsection{Общий угол и отношение площадей}

Если

\[
M\in AB,
\qquad
T\in AC,
\]

то треугольники $ABC$ и $AMT$
имеют общий угол $A$.

Формула площади

\[
S
=
\frac12 ab\sin\gamma
\]

даёт:

\[
S_{ABC}
=
\frac12
AB\cdot AC\cdot\sin A,
\]

\[
S_{AMT}
=
\frac12
AM\cdot AT\cdot\sin A.
\]

Делим первую формулу на вторую:

\[
\boxed{
\frac{S_{ABC}}{S_{AMT}}
=
\frac{AB\cdot AC}{AM\cdot AT}
}.
\]

Общий множитель

\[
\frac12\sin A
\]

сокращается.

\important{Универсальная идея.}
Если у двух фигур есть общий множитель в формулах,
запишите обе формулы и рассмотрите их отношение.

Часто лишние величины сократятся автоматически.

% =========================================================
% 8
% =========================================================

\clearpage
\section{8. Типичные ошибки и контроль решения}

\begin{itemize}

  \item
  \important{Высота к продолжению стороны.}

  Используйте в формуле площади длину стороны
  исходного треугольника и расстояние от вершины
  до прямой этой стороны.

  \item
  \important{Доверие рисунку.}

  Схема помогает рассуждать,
  однако равенство длин и углов
  должно следовать из условия или доказательства.

  \item
  \important{Слишком много неизвестных.}

  Новую переменную вводите тогда,
  когда через неё удаётся выразить
  несколько нужных элементов.

  \item
  \important{Недостаток независимых связей.}

  После введения переменных сразу оцените,
  хватает ли уравнений для их определения.

  \item
  \important{Тригонометрический контроль.}

  Значения синуса и косинуса
  должны лежать в допустимом диапазоне.

  \item
  \important{Вписанные углы.}

  Равенство углов связывайте
  с одной и той же хордой или дугой
  и учитывайте положение вершин на окружности.

  \item
  \important{Равнобедренный треугольник.}

  Сначала точно определите равные стороны,
  затем отметьте равные углы,
  лежащие напротив них.

  \item
  \important{Громоздкий путь.}

  Если теорема Пифагора порождает длинную систему,
  проверьте площадь, подобие, углы и окружность.

\end{itemize}

\subsection{Короткий чек-лист перед ответом}

\begin{enumerate}

  \item
  Искомая величина отмечена на рисунке.

  \item
  Все прямые углы и равные элементы
  имеют обоснование.

  \item
  Выбранная теорема применима
  к текущей фигуре.

  \item
  Каждая введённая переменная
  связана с условием.

  \item
  Арифметика и квадратные корни проверены.

  \item
  Длина положительна;
  тригонометрические значения
  находятся в допустимом диапазоне.

  \item
  Ответ соответствует единицам измерения
  и вопросу задачи.

\end{enumerate}

% =========================================================
% 9
% =========================================================

\clearpage
\section{9. Тренировочный блок --- Путь А}

\subsection{Средний уровень}

\begin{enumerate}[label=\textbf{\arabic*.}]

  \item
  В треугольнике две стороны равны $9$ и $6$.
  Высота к стороне $9$ равна $4$.
  Найдите высоту к стороне $6$.

  \item
  В $\triangle MJQ$ точка $T$
  лежит на стороне $MQ$, причём

  \[
  MT:TQ=5:3,
  \qquad
  S_{MJQ}=64.
  \]

  Найдите площади треугольников
  $MJT$ и $TJQ$.

  \item
  В остром $\triangle ABC$
  высоты $BD$ и $CE$
  пересекаются в точке $O$.

  Известно, что

  \[
  \angle A=62^\circ.
  \]

  Найдите

  \[
  \angle DOE.
  \]

\end{enumerate}

\subsection{Выше среднего}

\begin{enumerate}[label=\textbf{\arabic*.},resume]

  \item
  В равнобедренном $\triangle ABC$ выполнено

  \[
  AC=BC,
  \qquad
  AB=25.
  \]

  Высота $AH$ проведена к стороне $BC$,

  \[
  AH=20.
  \]

  Точка $H$ лежит на отрезке $BC$.
  Найдите $AC$.

  \item
  Точки $A,K,C$ лежат на одной прямой,
  $K$ находится между $A$ и $C$.

  Для треугольников $ABK$ и $KBC$
  высота из $B$ к прямой $AC$ общая.

  Известно:

  \[
  AK:KC=7:3,
  \]

  \[
  S_{ABK}=42.
  \]

  Найдите

  \[
  S_{KBC}.
  \]

  \item
  В $\triangle ABC$

  \[
  M\in AB,
  \qquad
  T\in AC.
  \]

  Известно:

  \[
  \frac{AM}{AB}
  =
  \frac25,
  \]

  \[
  \frac{AT}{AC}
  =
  \frac37,
  \]

  \[
  S_{AMT}=12.
  \]

  Найдите

  \[
  S_{ABC}.
  \]

  \item
  Точки $B,C,D,E$
  лежат на одной окружности,
  причём $BC$ --- диаметр.

  Точки $D$ и $E$
  расположены на одной полуокружности.

  Известно:

  \[
  \angle EBD=37^\circ.
  \]

  Найдите

  \[
  \angle ECD.
  \]

  \item
  Треугольник $BDC$ прямоугольный,

  \[
  \angle BDC=90^\circ,
  \qquad
  BC=10.
  \]

  Найдите радиус окружности,
  описанной около $\triangle BDC$.

  \item
  В остром $\triangle ABC$
  проведены высоты $BD$ и $CE$,
  где

  \[
  D\in AC,
  \qquad
  E\in AB.
  \]

  Известно:

  \[
  \angle ABC=68^\circ.
  \]

  Найдите

  \[
  \angle CDE.
  \]

\end{enumerate}

\subsection{Сложная задача}

\begin{enumerate}[label=\textbf{\arabic*.},resume]

  \item
  В тупоугольном равнобедренном
  $\triangle ABC$ выполнено

  \[
  AC=BC,
  \qquad
  AB=25.
  \]

  Высота $AH$ проведена к прямой $BC$
  и пересекает её за точкой $C$.

  Известно:

  \[
  AH=15.
  \]

  Найдите:

  \begin{enumerate}[label=\alph*)]

    \item
    длину $AC$;

    \item
    \[
    \sin\angle ACB.
    \]

  \end{enumerate}

\end{enumerate}

% =========================================================
% 10
% =========================================================

\clearpage
\section{10. Ответы}

\begin{table}[h!]
\centering

\caption{Ответы к тренировочному блоку}

\begin{tabularx}{\linewidth}{@{}cX@{}}

\toprule
\textbf{№} & \textbf{Ответ} \\
\midrule

1 &
$6$
\\

2 &
$S_{MJT}=40$,
$S_{TJQ}=24$
\\

3 &
$118^\circ$
\\

4 &
$\displaystyle \frac{125}{6}$
\\

5 &
$18$
\\

6 &
$70$
\\

7 &
$37^\circ$
\\

8 &
$5$
\\

9 &
$112^\circ$
\\

10 &
$\displaystyle AC=\frac{125}{8}$;
$\displaystyle \sin\angle ACB=\frac{24}{25}$
\\

\bottomrule

\end{tabularx}
\end{table}

\end{document}